\begin{center}
	\boxed{\text{In all this section, it is assumed that }
	\phaseAuthorizationList _{i} = \true
	}
\end{center}
We impose some general constraints onto \locAuthorizationIndex{}
and also ensure that \phaseAuthorizationList{}-phases have a certain shape.
Recall that \locAuthorizationIndex{} is meant to count and label authority tuples from the \texttt{authority\_list}
in order of appearance therein.
As such we must ensure it remain constant along the rows pertaining to a given authorization tuple.
Tuple labeling will start at $1$ and go all the way up to $\rlpTxnCommonColumnNumberOfAuthorizations$.
Furthermore tuple processing will always occupy $\locAuthorizationNumberOfRows$ rows,
regardless of the validity of the underlying tuple.

The constraints that follow pertain mostly to
\locAuthorizationIndex{} and enforce the following behaviour
on the \phaseAuthorizationList{} phase.
We wish to impose that these phases be structured as follows:
\begin{itemize}
	\item
		\numConst{1} ``transaction-row'' marks the beginning of
		the \phaseAuthorizationList{} phase, as it
		does for any of the phases;
		it is therefore expected;
	\item
		\numConst{1} follow-up ``computation-row'';
		again this is to be expected;
		on this row we compute the
		\rlp{}-prefix of the \rlp{}-ization of the
		\texttt{authority\_list};
	\item
		what follows are one or more batches of $\locAuthorizationNumberOfRows$ rows;
		every such batch pertains to a given authorization tuple;
		such batches should be characterized by a constant
		\locAuthorizationIndex{}, and this index should
		increment by one at the boundaries of said batches.
\end{itemize}
\label{rlp txn: phase constraints: authorization list: constraints: authorization lists must be nonempty}
The associated values of the \locAuthorizationIndex{} shorthand should look as follows
\[
	0, 0, ~
	\underbrace{1,    ~ \dots{} ~ , 1    } _{\displaystyle \locAuthorizationNumberOfRows \text{ terms}}, \quad
	\underbrace{2,    ~ \dots{} ~ , 2    } _{\displaystyle \locAuthorizationNumberOfRows \text{ terms}}, \quad \dots \quad, ~
	\underbrace{\ell, ~ \dots{} ~ , \ell } _{\displaystyle \locAuthorizationNumberOfRows \text{ terms}}.
\]
where we denote by $\ell$ the length of the \texttt{authorization\_list} i.e. $\ell \equiv \rlpTxnCommonColumnNumberOfAuthorizations$.

We achieve this by imposing the following constraints:
\begin{enumerate}
	\item\label{rlp txn: phase constraints: authorization list: constraints: setting values for the first few AUTH rows}
		\If $\locAuthorizationPrefix _{i} = 1$ \Then
		\[
			\left\{ \begin{array}{lclr}
				\locAuthorizationIndex  _{i}                                 & = & 0              \\
				\locAuthorizationIndex  _{i +                             1} & = & 1 \vspace{2mm} \\
				\phaseAuthorizationList _{i +                             1} & = & 1              \\
				\phaseAuthorizationList _{i +                            10} & = & 1              \\
				\phaseAuthorizationList _{i + \locAuthorizationNumberOfRows} & = & 1              \\
			\end{array} \right.
		\]
		\saNote{}
		We include ``look ahead checkpoint constraints'' of the form $\phaseAuthorizationList _{i + \relof} \equiv 1$
		in order to prevent premature exits from the present phase.
		The above ensures that the present phase lasts at least
		$1 + 1 + \locAuthorizationNumberOfRows$ rows.
		We will in the end ensure that the present phase lasts
		$1 + 1 + \ell \cdot \locAuthorizationNumberOfRows$ rows
		where $\ell$ is the length of the \texttt{authorization\_list} in the
		current ``set code transaction'' i.e. $\ell \equiv \rlpTxnCommonColumnNumberOfAuthorizations$.
	\item\label{rlp txn: phase constraints: authorization list: constraints: monotonicity of tuple index}
		\If $\phaseAuthorizationList _{i - 1} = \true$ \Then $\locAuthorizationIndex _{i} \in \{ \locAuthorizationIndex _{i - 1}, 1 + \locAuthorizationIndex _{i - 1} \}$
	\item \If $\locAuthorizationIndex _{i} \neq \rlpTxnCommonColumnNumberOfAuthorizations _{i}$ \et $\isCmp_{i} = 1$
		\Then
		\[
			\left\{ \begin{array}{lcl}
				\phaseAuthorizationList _{i + 1}                             & = & 1                               \\
				\phaseAuthorizationList _{i + \locAuthorizationNumberOfRows} & = & 1                               \\
				\locAuthorizationIndex  _{i + \locAuthorizationNumberOfRows} & = & 1 + \locAuthorizationIndex _{i} \\
			\end{array} \right.
		\]
		\saNote{}
		The above, along with the constraints setting the initial values for \locAuthorizationIndex{},
		see constraint~(\ref{rlp txn: phase constraints: authorization list: constraints: setting values for the first few AUTH rows}),
		the monotonicity constraint for \locAuthorizationIndex{},
		see constraint~(\ref{rlp txn: phase constraints: authorization list: constraints: monotonicity of tuple index}),
		enforce that past its two initial values \locAuthorizationIndex{} remains constant for $\locAuthorizationNumberOfRows$ consecutive rows
		and then grows by $1$.
		In other words we achieve the desired behaviour of
		diagram~(\ref{rlp txn: phase constraints: authorization list: constraints: authorization lists must be nonempty}).
	\item \If $\phaseAuthorizationList _{i + 1} = \false$ \Then $\locAuthorizationIndex _{i} = \rlpTxnCommonColumnNumberOfAuthorizations _{i}$
		% \item \If $\locAuthorizationFirstRow _{i} = 1$ \Then
		% 	\begin{enumerate}
		% 		\item $\locAuthorizationIndex _{i} = 1 + \locAuthorizationIndex _{i - 1}$
		% 		\item $\locAuthorizationIndex _{i + \locAuthorizationNumberOfRows - 1} = \locAuthorizationIndex _{i}$
		% 	\end{enumerate}
\end{enumerate}
